\section*{From finite spinors to a continuum Hamiltonian}
The preceding calculus becomes a concrete field construction when Fourier modes are used as packets. Each mode carries a four-component spinor and an explicit Hermitian matrix. The field is reconstructed through an isometric Fourier map. Exact norm preservation, completion, approximation, and control of the adjoint domain justify the passage from finite calculations to continuum analysis.

The setting is a free spinor on a fixed flat three-torus. For every real mass parameter, the resulting Hamiltonian is densely defined, self-adjoint, and closed. Its domain is exactly the first-order weighted Fourier square-summability domain. These assertions include the operator domain, not merely a formal differential expression.

The Hamiltonian and the modewise unitary evolution are first constructed separately. The following chapter, \emph{The Dirac Group and Its Generator}, connects them by proving strong continuity, equality of the full generator domain with the Hamiltonian domain, and the spectral-exponential identity. It also develops graph-domain and strong-resolvent convergence for the centered Fourier family and a quadratic spacing estimate under third-order Fourier regularity.

\begin{readerroadmap}[title={The role of this construction}]
Finite Clifford identities determine each mode. Hilbert-space completion reconstructs the field. Adjoint-domain equality determines the self-adjoint operator. The generator argument then identifies its evolution. These are distinct steps in one continuous mathematical construction.
\end{readerroadmap}

The operator construction, full generator-domain identification, centered graph and strong-resolvent limits, and third-order moment estimates have recorded Lean verification. The abstract spectral-functional-calculus identity remains a written deduction without a separate formal API linkage. The technical appendices identify the exact formal statements and evidence. None of these mathematical results by itself identifies an observed particle or supplies a dynamical spacetime geometry.

\clearpage
\section*{Reading routes and evidence}
\addcontentsline{toc}{section}{Reading routes and evidence}
Readers seeking the construction can begin with the conventions and Clifford matrices, follow the finite-packet estimates, and continue through coefficient completion to the adjoint calculation. Readers already comfortable with free Dirac matrices can begin at the full coefficient Hilbert space. The exact three-mode example gives a concrete state that can be followed through both continuum and centered evolution.

There are three kinds of statement in this volume. A \emph{checked correspondence} restates mathematics represented by named declarations in the retained formal source. An \emph{ordinary mathematical deduction} has its proof in the text but is not automatically assigned to a Lean declaration. An \emph{illustration} is an explicitly bounded Python evaluation of a formula. The figure record specifies its parameters and numerical checks.

The finite-packet propositions retain the identifiers PH01 to PH12 from the earlier traditional paper. This preserves their correspondence with its worked example and source notes. The later sections use numbered theorems for the completed-space and unbounded-operator results. A correspondence annotation identifies the checked source layer without asserting that the prose proof is a literal transcription of a tactic script.

\begin{longtable}{@{}L{.24\linewidth}L{.67\linewidth}@{}}
\toprule
Object & Meaning in this volume\\\midrule\endhead
Mode Hamiltonian & An everywhere-defined Hermitian matrix on a four-dimensional spinor space.\\
Coefficient evolution & A modewise unitary map on the complete coefficient Hilbert space.\\
Physical evolution & The conjugate of that map under the checked Fourier equivalence.\\
Continuum Hamiltonian & An unbounded partial linear operator with its domain carried explicitly.\\
Centered approximation & The multiplier with kinetic symbol $\sin(hp_j)/h$ on the same integer mode index set.\\
Generator identification & The theorem in the following chapter connecting the full-domain derivative and spectral exponential to the constructed Hamiltonian.\\
\bottomrule
\end{longtable}

The centered approximation is defined on continuum torus functions. A spatially sampled finite grid would require its own sampling, aliasing, and embedding conventions. Similarly, first-order regularity is expressed here by the precise Fourier condition that the sources prove. An identification with every possible weak-derivative or Sobolev-library convention would require a further equivalence theorem.

No finite count of successful numerical evaluations replaces the universal statements in the manuscript or their formal counterparts. Conversely, checking a mathematical model does not establish that the model describes observed matter. The physical setting here is a free spinor on a fixed flat periodic space, in natural units.

\clearpage
\section*{Resolution and finite dynamics}
\section{Packets and transfer blocks}
The useful starting point is a resolved operator rather than a preferred coordinate matrix. Let $V$ be a complex vector space with a finite family of projectors $(P_a)_{a\in A}$ satisfying
\[
P_aP_b=\delta_{ab}P_a,\qquad \sum_{a\in A}P_a=I.
\]
The component $P_av$ is the packet of $v$ selected by $a$. For a linear map $L:V\to V$, its transfer block from packet $a$ to packet $b$ is
\[
L_{ba}=P_bLP_a.
\]
Completeness gives $L=\sum_{a,b}L_{ba}$. Both indices matter. Keeping only $P_aLP_a$ would remove every transfer between distinct packets.

For two operators $L$ and $M$, insert the resolution between them:
\[
P_cMLP_a=\sum_{b\in A}(P_cMP_b)(P_bLP_a).
\]
This is block multiplication written without a choice of basis inside each packet. The intermediate index records the interface through which a composition passes. It is also the reason that a change of resolution must specify how fine packets assemble into coarse packets.

Suppose $P_A=\sum_{a\in C(A)}P_a$ and $P_B=\sum_{b\in C(B)}P_b$. Then
\[
P_BLP_A=\sum_{b\in C(B)}\sum_{a\in C(A)}P_bLP_a.
\]
The identity follows by distributing two finite sums. It says exactly which fine transfers a coarse observation retains. A scalar summary or a diagonal-only summary may retain less information and therefore needs its own loss statement.

These finite identities provide the notation used throughout the wider collection. The free Fourier model has an additional property: translation invariance makes its Hamiltonian diagonal in the mode label. Each diagonal block is still a nontrivial four-spinor matrix. The simplification belongs to this model, not to every operator that can be expressed in packets.

\subsection{What changes at completion}
The finite identity $\sum_aP_a=I$ becomes an approximation statement for an infinite orthogonal family. Finite partial sums must converge in the chosen topology. On the coefficient Hilbert space, canonical finite truncations converge in norm. For an unbounded Hamiltonian, approximating the state alone is insufficient to control its image. The graph approximation later in the volume controls both quantities at once.

This distinction explains the structure of the proof. Exact finite algebra supplies local formulas. Hilbert stability controls their extension to all square-summable fields. Weighted regularity specifies which fields support an unbounded operator action. Each step carries a different mathematical obligation.
